\section{Discussion}
\label{sec:discussion}

The results favour light over noise as the cause of early song in this population. Noise had no detectable effect once light was in the model, whereas the light effect was unchanged by the inclusion of noise. This does not show that noise is irrelevant to birdsong in general. It shows that, at the noise levels of ordinary residential streets, it does not move the start of the dawn chorus.

The saturation result matters for policy. Because most streets sit above the breakpoint, dimming a 10-lux street to 5 lux would be expected to change nothing. The schedule begins to recover only below about 3 lux, which is dimmer than current standards permit on most roads but typical of the part-night lighting schemes that several cities have adopted to save energy.

\paragraph{Limitations.} The study covers one species in one city, and the sites were chosen rather than sampled at random. We measured illuminance at recorder height and not at the roost, which probably understates the light that birds in dense vegetation experience. We did not measure fitness, so we cannot say whether singing earlier is costly.

\paragraph{Conclusion.} Street lighting advances the dawn chorus of urban blackbirds by up to half an hour, and traffic noise does not. The effect saturates at low light levels, so only substantial dimming is likely to reverse it.
